\documentclass[border=6pt]{standalone}
\usepackage[fontset=none]{ctex}
\usepackage{fontspec}
\setCJKmainfont{Songti SC}[BoldFont={Songti SC}]
\usepackage{amsmath,amssymb}
\DeclareMathOperator*{\argmax}{arg\,max}
\usepackage{xcolor}
\usepackage{tikz}
\usepackage{pgfplots}
\pgfplotsset{compat=1.18}
\usetikzlibrary{arrows.meta,positioning,shapes.geometric,calc,patterns}
\begin{document}
\begin{tikzpicture}[
  box/.style={draw,thick,rounded corners=2pt,minimum width=3.3cm,minimum height=1.05cm,align=center,fill=gray!8},
  lab/.style={font=\small,align=left,text width=4.6cm}
]
  \node[box] (a) at (0,0) {\textbf{事实层}\\ \small 政策改变了什么？};
  \node[box] (b) at (4.6,0) {\textbf{机制层}\\ \small 经何种约束/激励/均衡传导？};
  \node[box] (c) at (9.2,0) {\textbf{价值层}\\ \small 依何种福利标准评价？};
  \draw[->,very thick] (a) -- (b);
  \draw[->,very thick] (b) -- (c);
  \node[lab] at (0,-1.35) {\small 例：最低工资使低技能就业下降 2\%};
  \node[lab] at (4.6,-1.35) {\small 例：劳动需求弹性与覆盖范围决定效应大小};
  \node[lab] at (9.2,-1.35) {\small 例：用 $W=\sum\alpha_i u_i$ 还是罗尔斯标准？};
\end{tikzpicture}
\end{document}
